% ============================================================
%  Part XII — Usage
% ============================================================
\gpart{XIII}{Usage}

% ------------------------------------------------------------
\gsec{Impersonal es}

\gintro{German sentences need a subject even when nothing is really doing
anything. \textit{es} fills the gap, in three fairly distinct roles.}

\tcap{Weather and natural processes}

\begin{flattbl}{colspec={X[1,l] X[1,l] X[1,l] X[1,l]}, vline{3}={1-Z}{0.5pt,solid,Line}, column{3}={font=\small\sffamily\bfseries, fg=Ink}}
  es regnet  & it is raining  & es schneit & it is snowing \\
  es friert  & it is freezing & es donnert & it is thundering \\
  es ist kalt & it is cold    & es wird hell & it is getting light \\
\end{flattbl}

\tcap{es gibt — there is, there are}

\begin{extbl}{colspec={X[1,l] X[1.4,l]}}
  \textbf{Es gibt} einen Grund. & There is a reason. \\
  \textbf{Es gibt} viele Probleme. & There are many problems. \\
  \textbf{Gibt es} hier ein Café? & Is there a café here? \\
\end{extbl}

\begin{gnote}
\textit{es gibt} never changes for number and always takes the
\textbf{accusative}: \textit{es gibt \textbf{einen} Grund}, not \textit{ein
Grund}. Use \textit{es ist} or \textit{es sind} instead when you mean a specific
thing in a specific place: \textit{Es ist niemand zu Hause.}
\end{gnote}

\tcap{Verbs that describe how someone feels}

\begin{gtbl}{colspec={X[1.1,l] X[0.65,c] X[1.25,l]}}
  Expression & Case & Meaning \\
  es geht mir gut     & Dat. & I am well \\
  es gefällt mir      & Dat. & I like it \\
  es tut mir leid     & Dat. & I am sorry \\
  es fehlt mir        & Dat. & I miss it \\
  es ist mir kalt     & Dat. & I am cold \\
  es gibt             & Akk. & there is \\
\end{gtbl}

\begin{gnote}
\textit{es} gives up first position but stays in the clause:
\textit{\textbf{Es} geht mir gut} becomes \textit{Mir geht \textbf{es} gut}. In
\textit{Mir ist kalt} it drops out altogether.
\end{gnote}

% ------------------------------------------------------------
\gsec{Modal particles}

\gintro{Modal particles are short unstressed words that colour a sentence
without adding information. They are extremely common in speech, and leaving
them out is one of the things that makes correct German sound stiff.}

\begin{gtbl}{colspec={X[0.55,l] X[1.1,l] X[1.35,l]}}
  Particle & Effect & Example \\
  ja     & appeals to shared knowledge & Das ist \textbf{ja} klar. \eng{as we both know} \\
  doch   & urges, encourages           & Komm \textbf{doch} mit! \eng{do come} \\
  doch   & contradicts an assumption   & Das habe ich \textbf{doch} gesagt! \\
  mal    & softens a request           & Warte \textbf{mal}. \eng{hang on a sec} \\
  denn   & softens a question          & Was machst du \textbf{denn}? \eng{friendly} \\
  halt   & resignation                 & Das ist \textbf{halt} so. \eng{that's just how it is} \\
  eben   & same, slightly sharper      & Dann bleib \textbf{eben} hier. \\
  schon  & concedes a point            & Das stimmt \textbf{schon}, aber\dots \\
  wohl   & assumption                  & Er ist \textbf{wohl} krank. \eng{presumably} \\
  eigentlich & mild curiosity          & Wie alt bist du \textbf{eigentlich}? \\
  überhaupt & broadens the question    & Hast du \textbf{überhaupt} Zeit? \\
\end{gtbl}

\begin{gnote}
Particles sit in the middle field, never in first position, and are never
stressed. Several of these words also exist as ordinary adverbs with a full
meaning: \textit{schon} \emph{already}, \textit{eben} \emph{just now},
\textit{doch} \emph{yes it is}. Position and stress tell the two uses apart.
\end{gnote}

\tcap{doch as an answer}

\begin{extbl}{colspec={X[1,l] X[1.3,l]}}
  Du kommst nicht mit? — \textbf{Doch!} & You're not coming? — Yes I am! \\
\end{extbl}

\begin{gnote}
\textit{doch} is the word English lacks: as an answer it contradicts a negative
question or statement, where \textit{ja} would agree with the negative. In an
imperative it does the opposite job and simply urges.
\end{gnote}

% ------------------------------------------------------------
\gsec{False friends}

\gintro{These words look like English words and mean something else. They are
listed together because the errors they cause are hard to hear yourself
making.}

\begin{gtbl}{colspec={X[0.85,l] X[1,l] X[1.15,l]}}
  German & Actually means & English word is \\
  aktuell    & current, up to date & actually = eigentlich \\
  bekommen   & to receive          & to become = werden \\
  eventuell  & possibly            & eventually = schließlich \\
  sensibel   & sensitive           & sensible = vernünftig \\
  sympathisch & likeable           & sympathetic = mitfühlend \\
  spenden    & to donate           & to spend = ausgeben \\
  also       & so, therefore       & also = auch \\
  Gift       & poison              & gift = das Geschenk \\
  Chef       & boss                & chef = der Koch \\
  Rock       & skirt               & rock = der Stein \\
  Kind       & child               & kind = nett, freundlich \\
  Art        & kind, type          & art = die Kunst \\
  Handy      & mobile phone        & handy = praktisch \\
  brav       & well-behaved        & brave = mutig \\
  fast       & almost              & fast = schnell \\
  See        & lake \eng{der} / sea \eng{die} & to see = sehen \\
\end{gtbl}

\begin{gnote}
\textit{bekommen} is the costly one: \textit{Ich bekomme ein Bier} orders a
beer, it does not announce a transformation. The verb for \emph{to become} is
\textit{werden}.
\end{gnote}

\begin{gnote}
\textit{der See} is a lake and \textit{die See} is the sea — the same word with
two genders and two meanings. \textit{das Meer} is the safer choice for
\emph{sea}.
\end{gnote}
